% PAPER_A_METHOD_CONTRACT_CURRENT.json paper_a_method_20260807_v19 SHA-256 60e60f01d762ee89d788979ad9ef243db1887f6ef81f4d5bb8ae2e6f29b5af38
% !TeX program = xelatex
\documentclass[preprint,12pt]{elsarticle}

\usepackage{iftex}
\ifPDFTeX
  \PackageError{paper-a-outline}
    {This project requires XeLaTeX}
    {In Overleaf, open Menu and set Compiler to XeLaTeX.}
\fi

\usepackage[margin=2.3cm]{geometry}
\usepackage{amsmath,amssymb,mathtools}
\usepackage[strings]{underscore}
\usepackage{booktabs,longtable,tabularx,array,multirow}
\usepackage{enumitem}
\usepackage{graphicx}
\usepackage{xcolor}
\usepackage{hyperref}
\usepackage{fontspec}
\usepackage{xeCJK}

\setmainfont{TeX Gyre Termes}
\setCJKmainfont{FandolSong-Regular}
\setCJKsansfont{FandolHei-Regular}
\setCJKmonofont{FandolFang-Regular}
\hypersetup{
  hidelinks,
  pdftitle={可审计的自适应全景布局标注：从三轴测量到支持度感知前瞻政策}
}
\graphicspath{{figures/}}
\setcounter{secnumdepth}{3}
\setcounter{tocdepth}{2}

\begin{document}

\begin{frontmatter}
\title{可审计的自适应全景布局标注：\\
从三轴测量到支持度感知前瞻政策\\
\large 论文提纲与方法合同}
\author{匿名作者}

\begin{abstract}
本文研究全景布局标注中的模型辅助收益是否具有可测量的任务--工人异质性，
以及经 Calibration 支持的条件画像能否改善实际标注政策。
执行主线固定为
\(\mathrm{Pilot}\rightarrow\mathrm{P1}\rightarrow\mathrm{C1}\rightarrow\mathrm{C2\text{-}B}
\rightarrow\mathrm{C2\text{-}A\text{-}RP}\rightarrow\mathrm{T1}\rightarrow\mathrm{V1}\)。
C1/C2 将外部正确性、同行兼容性和结构可靠性分开测量，并以部分收缩控制小样本画像的不稳定性。
T1 采用 Manual/Semi-Auto 与普通/辅助风险的 $2\times2$ 设计；
V1 在共享工人池和对称容量约束下，前瞻比较 Strong Global 与 Support-aware Full。
协议显式区分工人失败、政策失败和外部系统事故，并以不可变证据、一次重跑和行政删失规则防止结果可见后的选择性归因。
V2 仅作为可选的外部支持与安全退化审计，不承担主要因果贡献。
本文不把尚未生成的数据、dry-run 工件或事后人工修订写成正式结果。
\end{abstract}

\begin{keyword}
全景布局标注 \sep 半自动标注 \sep 工人画像 \sep
自适应路由 \sep 前瞻政策试验 \sep 可审计协议
\end{keyword}
\end{frontmatter}

\tableofcontents
\clearpage

\part{简洁主文叙事}
% Overview layer: concise scientific narrative; detailed contracts follow.
\section{研究总览}

本文的统一科学问题是：全景布局标注中的模型辅助收益是否具有可测量的任务--工人异质性，
以及经 Calibration 支持的条件画像能否改善实际标注政策？统一逻辑为：

\begin{verbatim}
异质性存在
→ 异质性可以测量
→ 异质性可以由标注前任务信息激活
→ 异质性具有路由价值
\end{verbatim}

正式执行链为：

\begin{verbatim}
Pilot → P1 → C1 → C2-B → C2-A-RP → T1 → V1
\end{verbatim}

正文突出三项一级贡献：
\begin{enumerate}
  \item 分离外部正确性 \(Q^{GT}_u\)、同行兼容性 \(R^{peer}_u\) 和结构可靠性 \(F^{struct}_u\)；
  \item 以 common anchor、diverse bridge 和部分收缩建立双波 Calibration，避免发布不稳定的高维画像；
  \item 在相同容量、替补、追加和 GT-blind 聚合条件下，前瞻比较 Strong Global 与 Support-aware Full。
\end{enumerate}

P1 跨阶段预测、多峰共识和 GT 冲突解释是次级创新；反例库、失败归因、容量与调度审计是
可信度基础设施。V2 仅是可选的外部支持与安全退化审计。T1 与 V1 不是两个孤立实验：T1
检验模型辅助效应及其风险交互，V1 检验经 Calibration 支持的条件画像是否具有政策增量。


\clearpage
\part{完整方法与分析合同}
% Source: Paper_A_新版完整论文提纲_vFinal_Draft.md §1
\section{引言}

\subsection{研究背景}

全景房间布局标注需要工人同时处理：

\begin{itemize}
  \item 相机所在空间的识别；
  \item Scope/OOS 判断；
  \item 房间边界和角点几何；
  \item 模型初始化的识别与纠正；
  \item 遮挡、开放边界和邻接空间；
  \item 任务过程与完整性。
\end{itemize}

Semi-Auto 初始化可能降低操作负担，但也可能诱发：

\begin{itemize}
  \item blind trust；
  \item correction failure；
  \item over-correction；
  \item undercoverage；
  \item 结构性无效提交。
\end{itemize}

不同工人在一般任务、高风险任务和模型错误条件下的表现并不一定相同。单一的全局工人排名可能无法充分利用这种差异，但高维个人画像又容易受到小样本和支持不足影响。

\subsection{研究缺口}

现有方案通常存在以下不足：

\begin{enumerate}
  \item 把 PreScreen 仅视为准入，而没有检验其跨阶段预测价值；
  \item 将工人与同行一致性等同于外部正确性；
  \item 忽略结构性无效提交导致的选择性可计算问题；
  \item 在不同难度任务分发下直接按原始平均质量排名工人；
  \item 在小样本下无条件发布个体条件画像；
  \item 将路由简化为一个评分公式，而未定义 offer、容量、追加、聚合和终态；
  \item 只在 resolved 子集评价质量，忽视 unresolved 的选择性分母；
  \item 使用平衡试验分布替代真实生产任务分布。
\end{enumerate}

\subsection{方法路线}

本文采用：

\begin{verbatim}
高信息 P1
→ C1 基础测量
→ C2-B 共同桥接与风险确认
→ C2-A-RP 精度补测
→ T1 Semi 条件效应
→ Strong Global 对比 Support-aware Full V1
\end{verbatim}

并始终分离：

\begin{verbatim}
公开/纠错 GT 外部正确性
Geometry LOO 同行一致性
worker-caused structural failure
模型辅助风险
Calibration 支持范围
\end{verbatim}

\subsection{一级贡献}

\subsubsection{贡献一：三轴工人测量}

本文分别估计：

\begin{itemize}
  \item $Q^{GT}_u$：相对 operational GT 的外部正确性；
  \item $R^{peer}_u$：与独立同行的几何兼容性；
  \item $F^{struct}_u$：可归责于工人的结构失败倾向。
\end{itemize}

三者不压缩为单一总可靠度；P1 预测链和多峰共识作为次级效度证据。

\subsubsection{贡献二：双波 Calibration 与部分收缩}

C1 提供基础外部质量、同行结构和任务差异；C2-B 通过 common anchor 与 diverse
bridge 恢复横向可比性并确认风险差异，C2-A-RP 只补关键估计精度。
小样本工人使用部分收缩，不发布不稳定的高维画像。

\subsubsection{贡献三：支持度感知的前瞻政策试验}

V1 比较 Strong Global 与 Support-aware Full。Full 只能使用获得跨阶段支持的条件分量，
并在支持不足时局部或整体退化为 Global。两臂共享容量、替补、动态追加和 GT-blind 聚合条件；
完整工程合同放入方法附录。

P1 跨阶段预测、多峰共识、GT 冲突解释和反例库是次级创新或可信度基础设施，
不是替代三项核心贡献的主叙事。

\subsection{研究问题}

\subsubsection{RQ1：模型辅助效应}

Semi-Auto 是否在不降低交付质量的前提下降低 active time，其效应是否随任务风险和工人纠错能力变化？
确认性层级为 delivery-adjusted quality 非劣、active time 优势和
\texttt{mode} $\times$ \texttt{risk\_assist} 交互；blind trust、correction failure、
over-correction 和 Model Issue recognition 为机制性次级结果。

\subsubsection{RQ2：画像测量和跨阶段效度}

P1 诊断和 C1/C2 三轴工人状态是否具有跨阶段预测价值，哪些有限的 failure-family
component 能获得可复现支持？RQ2 以效应量、方向、支持范围和预测性能为主，
不要求每阶段分别达到传统显著性。

\subsubsection{RQ3：路由增量价值}

在相同工人池、容量、替补、追加和聚合规则下，Support-aware Full 是否相对
Strong Global 改善交付调整质量、非交付率或标注成本？V2 外部支持审计是可选扩展，
不作为第四个主要研究问题。

% Source: Paper_A_新版完整论文提纲_vFinal_Draft.md §2
\section{相关工作}

\subsection{全景布局估计与一维布局表示}

介绍 HorizonNet 的 ceiling--wall、floor--wall、wall--wall 一维表示，以及 HoHoNet 的 Latent Horizontal Feature。说明本文使用这些表示构造标注前模型风险，而非改进模型本身。

这一脉络还包括全景室内数据集、跨数据集布局评测和几何表示的域偏移问题。本文将模型输出
视为任务侧风险信号，并把模型不确定性与人工标注质量分开估计。

\subsection{模型辅助标注与 automation bias}

区分：

\begin{itemize}
  \item issue recognition；
  \item geometry correction；
  \item blind trust；
  \item over-correction。
\end{itemize}

说明工人勾选 Model Issue 不等于已成功纠正。

\subsection{众包工人建模与共识}

讨论：

\begin{itemize}
  \item 工人全局能力；
  \item 条件能力；
  \item 小样本收缩；
  \item disagreement；
  \item LOO independence；
  \item 多峰共识。
\end{itemize}

相关方法链还包括 worker--task difficulty 分解、任务难度建模、独立留一验证、经验贝叶斯
与层级收缩。本文沿用“工人能力、任务难度和观测噪声分离”的思想，但不把收缩后的能力
分数解释为外部正确性本身。

\subsection{分歧、多峰结构与外部真值}

众包研究表明，disagreement 可能来自任务本身的多解、工人策略差异或参考标准冲突；
因此多数票、单一 consensus 和 external GT 并不等价。本文先用 GT-blind 几何结构区分
unimodal、dominant-with-dissent 与 supported-multimodal，再将 GT 作为外部解释证据。

\subsection{自适应任务分配与动态冗余}

讨论：

\begin{itemize}
  \item worker selection；
  \item capacity；
  \item offer/acceptance；
  \item stopping；
  \item unresolved；
  \item policy evaluation。
\end{itemize}

动态冗余与自适应 worker selection 的研究还涉及 stopping、追加成本、候选池可用性、
容量约束、失败替补和政策级评价。本文把这些因素纳入前瞻状态机，并将 policy-caused
non-delivery 与 worker-caused structural failure 分开。

\subsection{人机协同、automation bias 与纠错}

人机标注工作通常区分 decision aid、automation bias、issue recognition、corrective
editing 与 over-correction。本文把模型辅助效果拆成 delivery-adjusted quality、结构
安全、active time 和机制性纠错结果，避免仅以 valid-only accuracy 评价协同系统。

\subsection{政策评价与数据中心的质量审计}

contextual routing、policy evaluation 和有限支持下的安全退化强调：路由规则必须在结果
可见前冻结，支持不足时回退到可解释基线，并以 ITT 与完整失败分母报告政策效果。数据中心
AI 与 annotation failure analysis 则要求记录 provenance、分母、审计证据和可复现版本；
本文的反例库、GT 冲突审查和结构失败事件均服务于这一可审计链条，而不是事后修改 GT 或
重新训练路由画像。

\subsection{可审计证据与信息时序}

强调：

\begin{itemize}
  \item 决策时只能读取当时已到达信息；
  \item 任务结果不能回填首次路由；
  \item 同一 submission 对不同 estimand 可以具有不同资格；
  \item 旧数据、修订 GT 和 worker state 必须版本化。
\end{itemize}

% PAPER_A_METHOD_CONTRACT_CURRENT.json paper_a_method_20260807_v19 SHA-256 60e60f01d762ee89d788979ad9ef243db1887f6ef81f4d5bb8ae2e6f29b5af38
\section{总体协议与数据生命周期}

\subsection{阶段职责}

\begin{longtable}{@{}p{0.08\textwidth}p{0.23\textwidth}p{0.31\textwidth}p{0.28\textwidth}@{}}
\toprule
阶段 & 核心职责 & 允许用途 & 禁止用途 \\
\midrule
\endhead
P1 & 准入、高信息诊断、预测因子 & admission、failure-family、P1$\rightarrow$后续预测 & 正式 Global、未经验证的路由画像 \\
C1 & 基础能力、风险发现、方差估计 & provisional worker state、C2 设计、功效模拟 & 最终政策、Main 结论 \\
C2 & 共同桥接、支持补齐、正式冻结 & Strong Global/Full、risk、support、fallback & 根据 Main 结果回调 \\
T1 & Semi 条件效应 & RQ1 & 修改 worker state \\
V1 & 路由前瞻验证 & RQ3 & 重新调权重和阈值 \\
V2 & 外部支持审计（可选） & 外部支持与安全退化 & 回流修改 V1 \\
\bottomrule
\end{longtable}

\subsection{冻结边界}

P1、C1 已发生的 admission、assignment 和 raw submission 不回改。

C2 后冻结：

\begin{verbatim}
reference registry
Global policy
Support-aware Full
risk_assist
risk_route
support thresholds
capacity
offer/replacement
dynamic redundancy
aggregation
terminal states
statistical analysis
\end{verbatim}

\subsection{C1 protocol deviations and authorized replacement}

原始 C1 assignment manifest、原始任务分配和 raw submission 保持不可变。现实执行中的
行政处置不覆盖原 manifest，而是通过独立的、带时间戳和 SHA 的
\texttt{authorized reassignment addendum} 追加：

\begin{itemize}
  \item W014 因 process-integrity failure 被行政排除；
  \item W034 与 W001 接手的替补任务只记录在授权替补 addendum 中；
  \item addendum 至少保存原 assignment manifest SHA、原任务/工人、替补任务/工人、
  授权人、授权时间、变更原因、addendum SHA 和生效时间。
\end{itemize}

因此，替补不是对原分配的静默重写。C1 assigned cohort 仍按原始 assignment 报告；
行政排除和替补任务分别进入 deviation、completion 和 analytic eligibility 字段，
并在结果中同时报告原始 assignment 与授权后实际执行状态。

对每一个基础任务 \(t\)，替补与追加的正式资格数定义为
\[
k_{t,\mathrm{eligible}}
=\#\{\text{unique eligible worker IDs on base task }t\},
\]
而不是 annotation rows 或 runtime IDs 的数量。替补工人不得此前提交过同一
\texttt{base\_task\_id}；该约束在 addendum 和最终任务级审计中同时核验。通过核验的
W034/W001 替补行不因其替补身份整体删除：只要满足对应资格，就可进入 GT、peer、LOO
和 structural-failure 分析；active-time 仅在该行同时满足 owner binding 与 active-time
资格时纳入。合法替补行也可进入主任务支持和终态聚合，并额外报告纳入与排除替补行的
profile sensitivity。

\subsection{信息时序}

\subsubsection{首次路由可读取}

\begin{verbatim}
公开任务 metadata
HoHoNet feature/output risk
Calibration worker state
availability
capacity
历史 process eligibility
\end{verbatim}

\subsubsection{响应到达后才可读取}

\begin{verbatim}
当前任务的合法提交
几何相似度
结构有效性
Scope state
冲突状态
已到达响应数
\end{verbatim}

\subsubsection{禁止回填首次路由}

\begin{verbatim}
当前任务 Difficulty
当前任务 Model Issue 票
最终 crowd consensus
最终 GT 评价
事后专家裁决
\end{verbatim}

\subsection{Estimand-specific inclusion}

同一 submission 分别保存：

\begin{verbatim}
eligible_for_active_time
eligible_for_GT_quality
eligible_for_LOO
eligible_for_structural_failure
eligible_for_semi_correction
eligible_for_predictive_validity
eligible_for_routing_feature
\end{verbatim}

不存在一个全局 \texttt{valid} 字段替代所有分析。

\subsection{失败归因、外部事故与重跑合同}

所有阶段将异常先按冻结证据规则区分为：

\begin{verbatim}
worker_caused_structural_failure
policy_caused_failure
external_system_failure
not_evaluable
\end{verbatim}

worker-caused structural failure 是可归责于提交的重复角点、非法 pair 数、自交、拓扑坍塌等；它进入对应 worker/condition 的结构机会分母。policy-caused failure 包括 candidate exhaustion、替补规则失败和容量耗尽，归责于原政策臂而非工人。external system failure 仅限有时间窗、服务端/平台/导出证据和影响范围记录支持的整站宕机、平台不可用或导出损坏；缺少该证据的异常为 not evaluable，不得事后改称系统事故。

C2 freeze 前必须锁定事故类别词表、证据要求、判定人/时间、每任务最多一次重跑、T1 配对重跑、V1 同臂同版本对称容量重跑和行政删失规则。分类、重跑或删失决定必须在查看受影响 T1 条件结果或 V1 政策结果前完成，保留 incident ID、原任务、重跑任务、证据 SHA 和 disposition。外部事故不降低 worker reliability，也不计为 policy failure；若无法按冻结规则重跑，则行政删失并按条件/臂报告原始、重跑、删失数量及原因。

\subsection{C1 assignment provenance、variable-k 与补充任务}

正式 C1 provenance 仅允许 \texttt{original\_assignment}、
\texttt{authorized\_replacement\_assignment}、\texttt{late\_entry\_calibration\_assignment}
与 \texttt{outside\_assignment\_submission}。前三类只有在 canonical identity、assignment binding
与对应 estimand eligibility 同时通过后才进入正式分析；outside 永远只进入 raw、exposure 和
process-integrity 审计。W014 永久行政排除。W034 的正式 capability set 是 original 与 17 条授权
任务的并集，W001 是 original 与 3 条授权任务的并集；W034 的 B-004/B-022 固定为 outside。

任务支持按 task--condition--estimand 定义：
\[
k_{t,c,e}=\#\{u:\text{worker }u\text{ 在 }(t,c,e)\text{ 上具有唯一 canonical eligible evidence}\}.
\]
GT、peer、LOO、structural、active-time、semi-correction、predictive-validity 与 routing-feature
可以具有不同的 \(k_{t,c,e}\)。原始 assignment 数只定义 target；授权替补恢复缺口，late entry
增加 pooled support，outside 和 duplicate revision 不增加正式支持。

W034 的 17 条授权任务只有在 owner-valid active-time sentinel 通过后、且任务开始时间不早于
sentinel validation time 时才贡献 timing evidence。sentinel manifest 必须绑定 worker、project、
runtime task、annotation、active-log SHA、reviewer 和 validation time。旧缺失日志不回填；sentinel
失败使相应 timing 分支 not-evaluable，不能伪造为零，也不能污染其他 capability 轴。

\subsection{Rolling enrollment 与 Stage 3 freeze}

Rolling enrollment 默认 \texttt{activated=false}，实际新增人数允许为 0。若在 Stage 3 前激活，
必须在结果不可见时冻结 \(N_{\max}\)、recruitment/P1/C1/C2 时间窗、P1/instruction/interface
版本、seed、workload template、task pool 与 exposure ledger。新人通过同版本 P1 后以独立 late-entry
manifest 追加，不替换 W014、不改写原 roster、不按结果选任务，并同时保留 original-only 与 pooled
profile。Stage 3 的启动由两套相互独立的 gate 控制，而不是一个等待所有后续政策工件的总门。
T1 仅要求 enrollment closed、全部 Calibration worker terminal、C1/C2-B/C2-A-RP、pooled profile
以及 T1 roster、task pool、randomization plan 和 SAP 的 SHA-bound freeze；因此 T1 不等待
Strong Global、Full 或 Validation roster。V1 在上述 Calibration/pooled 基础上另外要求
Strong Global、Full、Validation roster 和 V1 SAP 的 SHA-bound freeze。冻结后新增 worker 或修改
profile、threshold、weight、fallback、roster 均 fail closed。

T1 的正式 roles 为 \texttt{CALIBRATION\_ENROLLMENT\_CLOSED}、
\texttt{ALL\_CALIBRATION\_WORKERS\_TERMINAL}、\texttt{C1\_EVIDENCE\_FROZEN}、
\texttt{C2\_B\_FROZEN}、\texttt{C2\_A\_RP\_CLOSED}、\texttt{FINAL\_POOLED\_PROFILE\_FROZEN}、
\texttt{T1\_ROSTER\_FROZEN}、\texttt{T1\_TASK\_POOL\_FROZEN}、
\texttt{T1\_RANDOMIZATION\_PLAN\_FROZEN} 和 \texttt{T1\_SAP\_FROZEN}；V1 roles 在此基础上
替换为另外要求 \texttt{STRONG\_GLOBAL\_FROZEN}、\texttt{FULL\_POLICY\_FROZEN}、
\texttt{VALIDATION\_ROSTER\_FROZEN} 和 \texttt{V1\_SAP\_FROZEN} 的集合。每个 role 都必须
绑定对应冻结 artifact、method SHA 与递归依赖；仅增加角色字符串不能放行。

% Source: Paper_A_新版完整论文提纲_vFinal_Draft.md §4
\section{数据来源与单一 operational reference}

\subsection{GT 组成}

正式参考池由以下构成：

\begin{verbatim}
绝大多数：
公开数据集原始 GT

极少数：
经人工复核后纠错的 GT
\end{verbatim}

论文不得声称大规模重建 GT。

\subsection{Reference 状态}

\begin{verbatim}
reference_ready
reference_corrected
pending_adjudication
oos_geometry_not_applicable
\end{verbatim}

\subsubsection{\texttt{reference\_ready}}

直接使用公开原始 GT。

\subsubsection{\texttt{reference\_corrected}}

极少量明确问题经复核后冻结修订版本。

\subsubsection{\texttt{pending\_adjudication}}

单一 operational reference 尚未形成冻结裁决时，该任务不进入几何主分母。

\subsubsection{\texttt{oos\_geometry\_not\_applicable}}

当前相机房间协议无法形成合法单一布局，不进入普通几何评价。

\subsection{单一 operational reference}

正式表述为：

\begin{quote}
每张 in-scope、geometry-evaluable 任务，在冻结的 enclosed camera-room 协议下只保留一个 operational reference；无法唯一收口的任务进入 pending adjudication 或 OOS。
\end{quote}

不使用：

\begin{itemize}
  \item hard-multi；
  \item soft ambiguous reference；
  \item max-over-reference。
\end{itemize}

\subsection{GT 冲突两步审查}

\subsubsection{第一步：盲法初步裁决}

只查看：

\begin{verbatim}
原始全景图
公开 GT
冻结协议
\end{verbatim}

不查看工人身份、支持人数和工人等级。

\subsubsection{第二步：冲突解释}

初步裁决冻结后，再查看匿名提交，用于：

\begin{itemize}
  \item 解释 undercoverage；
  \item 解释 overextension；
  \item 识别协议误解；
  \item 建立反例类别。
\end{itemize}

裁决状态只允许：
\begin{verbatim}
retain_public_gt
correct_reference
oos_or_protocol_mismatch
unresolved_reference
\end{verbatim}

crowd cluster 必须先于 GT 判断形成；第一阶段只看图像、公开 GT 和冻结协议，
第二阶段才查看匿名 crowd submission。不得把 crowd majority 自动替换为 GT。

若第二步揭示明确遗漏并导致重新修订，标记：

\begin{verbatim}
submission_informed_reference_revision = true
\end{verbatim}

触发修订的原提交不使用该修订 GT 进入自身 primary GT-quality 评价，只进入敏感性或后续独立评价。
同时报告公开 GT、修订 GT 和排除修订任务的敏感性结果。

\subsection{数据切分}

硬要求：

\begin{itemize}
  \item C2、T1、V1 image-disjoint；
  \item 同一 image 的 Manual/Semi 在同一 T1 任务内；
  \item 同一 worker 不看同一 image 两种模式；
  \item T1/V1 尽量 building-disjoint；
  \item ZInD 仅用于 V2。
\end{itemize}

% PAPER_A_METHOD_CONTRACT_CURRENT.json paper_a_method_20260807_v19 SHA-256 60e60f01d762ee89d788979ad9ef243db1887f6ef81f4d5bb8ae2e6f29b5af38
\section{同行结构、GT 解释与三状态观测}

\subsection{四类证据}

\begin{verbatim}
Scope/OOS
Difficulty
Model Issue
Geometry
\end{verbatim}

四类证据进入不同的资格、风险和解释 gate。GT 只用于外部正确性评价与冲突解释，
不能反向定义 crowd structure，也不能把多数提交自动替换为 GT。

\subsection{GT-blind crowd structure}

在查看 GT 解释之前，先对所有合法提交按冻结几何相似度形成兼容簇。任务级结构标签为：

\begin{verbatim}
unimodal
dominant_with_dissent
supported_multimodal
insufficient_or_incompatible
\end{verbatim}

定义如下：
\begin{itemize}
  \item \texttt{unimodal}：所有合法提交形成同一兼容簇；
  \item \texttt{dominant\_with\_dissent}：存在唯一主簇，第二簇支持不超过 1；
  \item \texttt{supported\_multimodal}：至少两个各有两名或以上支持的兼容簇；
  \item \texttt{insufficient\_or\_incompatible}：支持不足或几何不可比较。
\end{itemize}

因此，4:1 和 3:1:1 通常属于 \texttt{dominant\_with\_dissent}；真正的 3:2 属于
\texttt{supported\_multimodal}，前提是两个簇各自内部一致。若 3:2 中两人簇内部也不一致，
则标为 \texttt{dominant\_with\_dissent} 或 \texttt{insufficient\_or\_incompatible}。

\subsection{三类共识的职责}

\begin{itemize}
  \item 任务描述共识：描述场景或问题是否被多名工人主张；
  \item 几何同行结构：形成兼容簇、medoid、margin 和多峰状态；
  \item 停止共识：决定是否追加和是否形成 resolved 输出。
\end{itemize}

三者不共用一个多数阈值。3:2 不生成唯一 crowd truth；所有工人仍可计算
\texttt{R\_peer\_task}，并按 task 等权汇总为 \texttt{R\_peer\_all}，稳定敏感性使用
\texttt{R\_peer\_stable}；两个簇分别与 GT 比较，少数簇不自动定义为错误。

\subsection{三状态任务观测及 \texttt{not-evaluable} 状态}

对 task--tag 保存三种任务观测，并将 \texttt{not-evaluable} 作为独立的分析状态：

\begin{verbatim}
positive assertion
explicit negative
unasserted
not evaluable
\end{verbatim}

以及 \(a=\) positive count、\(e=\) explicit negative count 和
\(u=\) unasserted count。若 \(a\geq2\) 且 \(e\geq2\)，记为显式冲突，
不转成普通 positive/negative truth。\texttt{not-evaluable} 不属于第四种任务标签，
而表示证据不足、协议不适用或归因未决时的分析排除状态；它不进入相应标签的分母，
也不被静默编码为 negative。

\subsection{Scope/OOS、undercoverage 与 Model Issue}

OOS 表示任务是否适合当前协议；undercoverage 表示任务可标但工人未完整覆盖相机所在空间，
二者不得混淆。Model Issue 进一步区分 \texttt{issue recognition}、\texttt{geometry correction}、
\texttt{blind trust} 和 \texttt{over-correction}。反馈栏耗时不单独扣除，完整任务 active time
仍是 T1 的真实使用成本。

\subsection{任务级同行统计与 crowd support}

pairwise edge 只用于构造每个 worker--task 的同行兼容度，不作为额外独立样本。论文先计算
\(R^{peer}_{t,u}\)，再在 worker 层对 task 等权汇总，避免高 \(k\) 任务因边数较多获得更高统计权重。
每个任务报告最大簇和次大簇的 absolute support、share、cluster count，以及当前字段合同的
\texttt{cluster\_margin\_all} 与 \texttt{cluster\_margin\_top2}；不再生成名为
\texttt{normalized\_cluster\_margin} 的字段。supported multimodal 任务保留 peer 描述与 GT
对齐审计，但不进入 stable-peer tie-break。

% Paper A: P1 is a secondary predictive contribution and an audit layer.
\section{P1：高信息准入、诊断与预测}

\subsection{P1 的定位}

P1 不仅用于准入，也用于形成可跨阶段检验的 failure-family 假设；但 P1 不承担正式
Global 分数。完整诊断词表保留在审计层：

\begin{verbatim}
undercoverage
adjacent-space overextension
corner instability
topology / over-parsing
blind trust
correction failure
scope misunderstanding
process-integrity risk
\end{verbatim}

其中 blind trust、correction failure 主要服务 T1 机制分析；Scope 是资格和协议审计；
process-integrity 是安全门。它们不进入 V1 条件加权。

\subsection{跨阶段预测链}

\begin{verbatim}
P1 → C1
P1 → C2-B
P1 → T1
\end{verbatim}

报告 Spearman/Kendall、worker bootstrap CI、方向一致性、discrepancy workers、support
和 range restriction。RQ2 关注效应量、方向、支持范围和预测性能；不要求 P1、C1、C2
每阶段分别达到传统显著性。

\subsection{正式可路由 failure family}

正式 Full 只考虑以下三个可由标注前任务特征激活的 family：

\begin{enumerate}
  \item undercoverage / underextension；
  \item adjacent-space overextension；
  \item corner / topology instability。
\end{enumerate}

component 必须同时通过 C1 predictive validation 与 C2-B confirmation，方向符合预设假设，
工人和任务 support 足够，个体效应经过收缩，并且在任务侧可激活。证据不足时调整自动趋近零，
C2-A-RP 只补已定义 component 的精度，不发现新 family。

\subsection{P1 retrospective integrity amendment}

\begin{itemize}
  \item parent-derived 或其他 non-independent P1 submission 不进入 capability 或 timing 主证据；
  \item 这些记录可保留为 process-integrity 负证据，但不回改已经完成的 admission；
  \item P1 预测分析只使用 independent 且 estimand-eligible 的证据；
  \item 未经 C1+C2-B 双门验证的 component 保持 diagnostic-only。
\end{itemize}

\subsection{轻量反例库}

反例库由以下来源触发：

\begin{verbatim}
GT conflict
supported multimodality
structural failure
blind trust / correction failure
undercoverage / overextension
V1 unresolved
\end{verbatim}

论文合同要求候选至少保存 \texttt{candidate\_id}、来源、触发规则版本、证据、分母定义和
输入 SHA；候选不自动表示 failure、GT 错误或工人归责。若任一必需 provenance、分母
定义或 evidence SHA 缺失，候选只能保留在完整性审计中，不进入正式失败率、富集率或
代表性反例分母。
独立裁决类别为：

\begin{verbatim}
confirmed_worker_failure
confirmed_ambiguity
reference_or_protocol_issue
system_or_policy_issue
not_counterexample
unresolved
\end{verbatim}

本轮只选择少量代表案例用于正文图、结果解释、讨论和附录；完整 frozen challenge/regression
bank 是未来可选用途，不是本轮硬交付。反例库不得用于当前画像、C2 阈值、Full 权重或总体
prevalence，也不得自动转为修订 GT。

% PAPER_A_METHOD_CONTRACT_CURRENT.json paper_a_method_20260810_v20 SHA-256 085fd9ec0f14a93c986e9f7dda7883173d9ec64959188139a25012d0e6d36b9d
% Paper A: C1 separates computability from formal analysis eligibility.
\section{C1：三轴工人测量与设计参数}

\subsection{C1 目标}

C1 提供基础外部质量、同行结构、任务差异和结构失败机会，并为 C2-B、T1 和 V1
提供预注册设计参数。C1 不把未通过正式资格门的可计算分数升级为 Global 支持。

\subsection{外部正确性与资格}

\[
Q^{GT}_{t,u}=IoU(A_{t,u},G_t).
\]

将原来的单一 \texttt{quality\_evaluable} 拆分为：

\begin{verbatim}
gt_score_computable
gt_primary_analysis_eligible
\end{verbatim}

前者只表示 IoU 可以计算；后者还必须通过 process、independence、Scope、reference、
Manual condition 和 geometry computability。结果表不得把前者称为正式 Global 支持。

\subsection{Task-adjusted GT quality 与经验贝叶斯收缩}

在 C1 阶段，先使用不含 stage effect 的 provisional 模型：

\[
Q^{GT}_{t,u}=\mu+\alpha_u+\delta_t+\epsilon_{t,u}.
\]

C1 与 C2 合并后，才估计含阶段效应的最终模型：

\[
Q^{GT}_{t,u,s}=\mu+\alpha_u+\delta_t+\eta_s+\epsilon_{t,u,s}.
\]

在 task-adjusted fixed-effect 估计之外，使用正常--正常经验贝叶斯收缩，报告：

\begin{verbatim}
Q_GT_task_adjusted_FE
Q_GT_EB
Q_GT_EB_interval
Q_GT_LCB
\end{verbatim}

正式 Global 默认使用收缩后的 \(Q^{GT,EB}_u\)；LCB 是最低质量安全门、
不确定性标记和保守敏感性排序，不再是唯一的正式排序值。

\subsection{同行兼容性的双指标}

两条几何通道先分别计算，再定义同行相似度：

\[
S(A_{t,u},A_{t,v})
=\min\left\{S_{\mathrm{boundary}}(A_{t,u},A_{t,v}),
S_{\mathrm{wall}}(A_{t,u},A_{t,v})\right\}.
\]

这一定义对应“任一关键几何通道明显不兼容时，不用另一通道的高分掩盖”的合同。

主要连续指标为：

\[
R^{peer}_{t,u}=\operatorname{median}_{v\neq u}
S(A_{t,u},A_{t,v}),
\]

回答工人与同行总体的兼容程度。高置信参考指标为：

\[
R^{peer\_all}_u=\operatorname{median}\{S(A_{t,u},A_{t,v}):
v\neq u,\ v\text{ eligible}\},
\]

\[
R^{peer\_nonmultimodal}_u=\operatorname{median}\{S(A_{t,u},A_{t,v}):
v\neq u,\ v\text{ eligible},\
t\notin\texttt{supported\_multimodal}\}.
\]

\texttt{R\_peer\_all} 用于描述性报告；Global 的 tie-break 与 stable-compatibility
判断只能使用 \texttt{R\_peer\_stable}。支持多峰任务只进入 discrepancy 与
敏感性分析，不进入 Global 的稳定兼容性证据。

\[
R^{LOO}_{t,u}=S(A_{t,u},M_t^{(-u)}),
\]

其中 \(M_t^{(-u)}\) 只有在排除工人后存在唯一稳定主簇时才定义。论文明确区分：

\begin{verbatim}
R_peer_stable       主要同行兼容指标
R_LOO_medoid        高置信参考、聚合稳定性与 tie-break
R_LOO_strict       保守敏感性分析
\end{verbatim}

三者不得机械平均，也不得用 LOO 填充 GT quality。

\subsection{结构失败与收缩}

原始结构失败比例保留为：

\[
F^{struct}_{u,raw}=\frac{n_{\mathrm{worker\text{-}caused\ failure}}}
{n_{\mathrm{structural\text{-}evaluable\ opportunities}}}.
\]

同时报告：

\begin{verbatim}
F_struct_raw
F_struct_EB
F_struct_interval
recurrent_serious_failure_flag
\end{verbatim}

external system failure、parser/export failure、reference failure、OOS 和未知归因不进入
worker-caused 分母。正式安全门由独立 policy manifest 冻结，本文不写死任意固定数值阈值。

\subsection{三轴基础状态}

\[
Q^{GT}_u,\qquad R^{peer}_u,\qquad F^{struct}_u
\]

分别代表外部正确性、同行兼容性和结构可靠性。三轴分开估计；同行指标只用于 tie-break、
discrepancy worker、聚合稳定性和敏感性分析，不作为外部质量的替代物。

\subsection{Estimand-specific support 与 profile 版本}

每个 worker profile 必须携带 GT、peer、LOO、structural 与 timing 的独立 task support、deficit
和 excess；不得用一个全局 (k=5) 或 annotation row 数替代。W034 同时物化 original-only 与
original+authorized profile，W001 对应 original+authorized profile；rolling enrollment 激活时再增加
original-cohort 与 pooled-cohort 版本。所有 profile 都绑定 assignment、canonical evidence、row
eligibility 和 estimator manifest 的 SHA。Active-time sentinel 仅约束 timing validation；其缺失使 timing
not-evaluable，但不得阻断几何三轴、C1-A snapshot 或 C2-B roster。Label Studio lead time 不作为逐行
fallback，也不与 active time 混合；若统一计算，只能报告为独立的 post-hoc exploratory elapsed-time。

\subsection{三轴用途级资格合同}

正式 worker profile 分别报告 \(Q_{GT}\)、\(R_{peer}\)、\(F_{struct}\)、
\(R_{LOO,medoid}\) 与 \(R_{LOO,strict}\) 的状态。C2-B roster 只要求前三个基础轴可评价；
strict/medoid LOO 不构成全局 roster 硬门。同行轴先在 worker--task 内取 pairwise similarity
中位数，再在 worker 层对 task 等权取中位数，不直接汇总全部 pairwise edge。

\paragraph{Contract v10 clarification.}
The authoritative C1-only model is \(Q^{GT}_{t,u}=\mu+\alpha_u+\delta_t+\epsilon_{t,u}\), where \(\delta_t\) is a task fixed effect and no stage effect is fit. The C1+C2 final model is \(Q^{GT}_{t,u,s}=\mu+\alpha_u+\eta_s+b_{\mathrm{building}(t)}+b_{\mathrm{task}(t)}+\epsilon_{t,u,s}\); it is reported as \texttt{not\_identifiable} without a frozen cross-stage anchor or equivalent supported structure. Define \(R^{peer}_{t,u}=\operatorname{median}_{v\ne u}S(A_{t,u},A_{t,v})\), \(R^{peer\_all}_u=\operatorname{median}_{t\in\mathcal T_u}R^{peer}_{t,u}\), and \(R^{peer\_stable}_u=\operatorname{median}_{t\in\mathcal T_u\setminus\mathrm{supported\_multimodal}}R^{peer}_{t,u}\). These definitions replace any earlier all-edge aggregation wording in this section.

\paragraph{Task-level Scope closure.}
Scope is a task-level reference decision rather than a worker axis. The complete 75-task researcher first review is compared with the cleaned canonical worker group direction as descriptive evidence (at least three responses and a unique relative majority); formal submission independence is not required for this Scope triage and remains required separately for Q\_GT, peer, LOO, and structural estimands. Concordant directions are frozen directly; disagreement, a tie, or lower support receives a second protocol review. The terminal disposition is \texttt{in\_scope}, \texttt{oos}, or \texttt{unresolved}. Only \texttt{in\_scope} tasks enter Q\_GT, peer, LOO, and structural opportunities. A worker's OOS response on a final in-scope task does not itself discard an otherwise eligible geometry submission. Terminal \texttt{oos} and \texttt{unresolved} tasks remain in audit evidence but never create a worker-caused structural failure.

% Source: Paper_A_新版完整论文提纲_vFinal_Draft.md §8
\section{HoHoNet/HorizonNet 任务风险}

\subsection{三个变量}

\begin{verbatim}
d_model_feat
g_model_struct
d_cal_support
\end{verbatim}

\subsection{\texttt{d\_model\_feat}}

基于 HoHoNet LHFeat：

\[
d^{feat}_t
=
kNNDist\left(
\phi(H_t),
\{\phi(H_i):i\in B_{\text{train-ref}}\}
\right).
\]

处理：

\begin{itemize}
  \item checkpoint 固定；
  \item preprocess 固定；
  \item PCA/whitening 只在参考库拟合；
  \item 处理 circular shift 和 seam；
  \item 保存 global 与 local-max 距离。
\end{itemize}

\subsection{\texttt{g\_model\_struct}}

基于一维布局输出：

\begin{verbatim}
ceiling/floor curve
wall-wall peaks
corner/pair count
topology
seam stability
postprocess validity
\end{verbatim}

保存可解释 flags，不只保留单一风险分。

\subsection{\texttt{d\_cal\_support}}

\subsubsection{\texorpdfstring{\(d_{\text{cal}}^A\)}{dcal-A}}

只基于 C1，用于设计 C2。

\subsubsection{\texorpdfstring{\(d_{\text{cal}}^F\)}{dcal-F}}

C2 后计算，用于 V1 激活和 fallback。

不得用最终版本反向解释 C2 选择。

\subsection{两类风险}

\subsubsection{\texttt{risk\_assist}}

服务 T1，预测模型初始化与纠正风险。

\subsubsection{\texttt{risk\_route}}

服务 V1，预测 Manual 外部质量下降、LOO disagreement、结构失败或 worker ranking instability。

C1 用于候选筛选，C2-B 用于共同确认，C2-A 只补精度，不重新搜索新风险。

% PAPER_A_METHOD_CONTRACT_CURRENT.json paper_a_method_20260807_v19 SHA-256 60e60f01d762ee89d788979ad9ef243db1887f6ef81f4d5bb8ae2e6f29b5af38
\section{C2：共同桥接、层级收缩与精度补齐}

\subsection{C2-B 的有限候选设计}

C2 只保留三个正式候选：

\begin{verbatim}
D8  = 4 common anchors + 4 diverse bridges
D10 = 6 common anchors + 4 diverse bridges
D12 = 6 common anchors + 6 diverse bridges
\end{verbatim}

选择满足以下要求的最小设计：
\begin{itemize}
  \item worker--task 图连通；
  \item common anchor 足够且 ordinary/stress 均有支持；
  \item task/building 多样性达标；
  \item Global 排名达到预定稳定性；
  \item 风险斜率达到预定精度；
  \item 总预算可接受。
\end{itemize}

政策差异率只作为 V1 可识别性诊断，不作为 C2 设计的主要优化目标。

\subsection{共同 anchor 与 diverse bridge}

common anchor 用于横向可比和波次校正；diverse bridge 用于恢复 task/building 外推能力、
ordinary/stress 支持和 worker--task 图连通性。候选设计、模拟输入和最终选择必须在
Calibration 结果可见的边界内预先冻结，不从反例库或 V1 outcome 反推。

\subsection{风险韧性层级模型}

\[
Q^{GT}_{u,t}=\alpha_u+\gamma I_t^{risk}
+b_uI_t^{risk}+\eta_{\mathrm{stage}}+r_t+\epsilon_{u,t},
\qquad b_u\sim N(0,\tau_b^2).
\]

输出收缩后的 risk slope、interval、leave-one-task/block-out stability 和 routing eligibility。
不确定性大时调整趋近零，不发布不稳定的高维 worker\(\times\)scene 画像。

\subsection{C2-A-RP：精度自适应补测}

C2-A-RP 的正式派发目标只有 \texttt{risk\_slope\_precision}。top-$k$ 边界与 component
eligibility 只保留为诊断字段，不得作为正式派发目标：

\begin{itemize}
  \item 每次补测固定为 1 ordinary + 1 stress；
  \item 每人最多五个 block；
  \item 最多增加 10 张任务；
  \item 只更新已声明的条件分量，不搜索新风险族。
\end{itemize}

正式区间统一使用冻结的 \texttt{normal\_95\_max\_unified\_slope\_sd} 规则和 95\% 正态区间；
目标宽度只从 C2-B 结果可见前冻结的 threshold manifest 读取，并绑定其 SHA，不使用硬编码值
或简单支持量缩放公式。选择依据是预期减少决策不确定性，而不是制造政策差异或提高显著性。
每次补测保存 target component、gap reason、候选池、历史排除、随机种子、declared/projected
support 和实际 observed support。每个 block 固定为 1 ordinary + 1 stress，最多五个 block，
因此每人正式任务数只能是 0、2、4、6、8 或 10；每个 block 完成后，从实际 canonical、valid 且
risk-slope-estimand-eligible 的提交重新拟合同一模型，更新 estimate、SE、CI、half-width 和
ordinary/stress support。达到目标即停止；五块后仍未达标时 risk adjustment=0 并记录
\texttt{fallback\_action=STRONG\_GLOBAL}。target-met、fallback 与 not-evaluable 互斥，
所有 plan worker 必须有终态；0-task closeout 只有在所有 worker 已终态时才合法。

若 D8、D10、D12 均未同时通过图连通、Global 稳定性或风险精度门，不扩展候选设计，
不追加未经预声明的新 component；所有未通过的 component adjustment 置零。若 Strong
Global 仍可冻结，则保留 Global/T1 并不启动 Full/V1；若 Global 本身不稳定，则不启动
V1。该规则优先于追求政策差异或补救性扩大样本。

\subsection{C2-S：可选诊断}

C2-S 仅可作为 Semi correction support 的可选诊断，不是正式主线阶段，不进入 Manual Global
或 V1 的主要 Full 设计。

\subsection{C2 冻结门}

正式 C2-A-RP 输入缺失、formal goal 非法、任一 SHA 不一致或 C2-B roster 不完整时，
只生成带 reason code 与 hashes 的 blocked summary，以非零状态结束，不生成
\texttt{launch\_ready=true} assignment。closeout 同时保留 declared/projected support 与
observed support；计划值不得被实际值静默覆盖。C2-B closeout 只作为审计证据，不能直接
替代 Full Policy。

T1 只要求 Calibration、pooled 与 T1-specific 的十项 freeze roles；V1 才另外要求 Strong Global、
Full、Validation roster 与 V1 SAP。T1 的 2$\times$2 条件试验不以 Global 排名替代；启动 V1
前则必须确认 C2-B 设计、C2-A-RP、Strong Global、Full 和 Validation roster 的各自 artifact
均已按 V1 gate 冻结。

% PAPER_A_METHOD_CONTRACT_CURRENT.json paper_a_method_20260730_v5 SHA-256 bde2e7e20cb00fa4f67b377112fe6534e27e7938c34fb4f63b7987fd3c142e2b
% Paper A: Strong Global is the EB external-quality baseline.
\section{Strong Global 与 Support-aware Full}

\subsection{Global 候选审计}

Global-LOO、Global-GT 和 Global-GT-gated-LOO 仍可在 cross-fitted replay 中比较，
报告排名重叠、discrepancy worker 和稳定性；这些只是审计和敏感性分析，不能替代正式
Strong Global，也不能由 V1 outcome 选择。

\subsection{Strong Global}

\subsubsection{Eligibility}

\begin{verbatim}
process valid
independence valid
GT support sufficient
building/task coverage sufficient
no recurrent serious structural failure
Q_GT_EB meets the production quality floor
\end{verbatim}

\subsubsection{正式分数与安全门}

\[
S^G_u=z\!\left(\widehat Q^{GT,EB}_u\right).
\]

其中 \(z(\cdot)\) 的均值和标准差只在冻结的、行政上具备资格且可估计
\(Q^{GT}\) 的 capability cohort 上计算，并在质量 floor 与 structural gate 之前完成。
通过质量 floor 或 structural gate 后不重新标准化；Strong Global 的正式排序不能借由
筛选后的资格变化改写尺度。

正式排序使用收缩后的外部正确性；LCB 只作为最低质量安全门、不确定性标记和保守敏感性
排序。同行指标不进入正式分数，仅用于 tie-break、discrepancy worker、聚合稳定性和敏感性。

\subsection{Support-aware Full}

\[
S^F_{u,t}=S^G_u
+I_{\mathrm{risk},t}B_u
+I_{f,t}H_{u,f}.
\]

其中 \(B_u\) 表示 ordinary/stress 风险韧性，\(H_{u,f}\) 表示一个已激活 failure-family
的条件能力。每任务最多激活一个 family，总调整有上限；证据有限时向群体均值收缩，
component 无支持时该项为零，任务超出 Calibration 支持范围时整体 fallback Strong Global。

\subsection{可路由 family 与激活}

正式 Full 只允许：

\begin{enumerate}
  \item undercoverage / underextension；
  \item adjacent-space overextension；
  \item corner / topology instability。
\end{enumerate}

\texttt{blind trust} 与 \texttt{correction failure} 仅作为 T1 机制结果；Scope 和
process-integrity 不进入 V1 条件加权。任务侧 family 必须唯一激活；激活含糊则关闭该项。

\subsection{支持、权重与 fallback}

跨阶段 component support 是预声明的方法合同。对每一个预定义 component，使用任务级
Calibration outcome 的 stage-stacked 回归或层级模型，包含 P1 worker component、stage、
P1\(\times\)stage interaction，并按 task 与 worker 聚类或设置相应随机效应。模型输出
方向、stage interaction、worker/task support 与收缩估计，并在 Main outcome 可见前冻结。
若 C2 出现清晰的方向反转、支持不足或不稳定收缩，该 component 不进入 Full。

只有当上述模型工件和独立 policy manifest 均有效时，component 才能通过 P1 integrity、
C1 predictive validation 和 C2-B confirmation，并具备足够 worker/task support、方向稳定、
收缩估计和标注前可激活性。权重与总调整上限不能为了增加政策差异而改变。

\begin{verbatim}
component unsupported        -> component adjustment = 0
family activation ambiguous  -> component adjustment = 0
d_cal_F out of support       -> overall fallback Strong Global
profile version incompatible -> overall fallback Strong Global
ranking unstable             -> overall fallback Strong Global
\end{verbatim}

整体 fallback 时，Full 的排序逐项等于 Strong Global。C2 closeout 后仅报告 activation、
fallback、推荐不同率、supported candidate count 和 capacity 后差异，不能用 V1 结果放宽门槛。

\subsection{Stage 3 formal manifest contract}

Strong Global 的 \(z\)-score cohort 在 quality/structural gate 前冻结为 administratively eligible 且
\(Q^{GT,EB}\) estimable 的 worker；cohort 少于两人、标准差为零、字段缺失或非有限值时 fail closed。
LCB 只作 quality floor、区间和 sensitivity，不参与正式主排序。Full manifest 必须显式绑定
component whitelist、worker/task support、weight、activation threshold/margin、symmetric cap、
component interval、profile version、fallback、formal minimum worker/cluster rule 与全部输入 SHA。
任何字段缺失、版本冲突、支持不足或 ranking endpoint 不稳定时，Full 整体排序精确回到 Strong Global。
V1 freeze manifest 还必须消费开放且 SHA 匹配的 Stage 3 gate；GT-blind aggregation 不读取
\(Q^{GT}\)、\(S^G\)、\(S^F\) 或 worker quality 作 medoid tie-break。

\subsection{静态画像与运行时调度分离}

Strong Global 的唯一静态顺序为
\(S_G\rightarrow R_{\mathrm{peer,stable}}\rightarrow R_{\mathrm{LOO,medoid}}\rightarrow\)
frozen random。availability 与 capacity 只在运行时 scheduler 过滤候选，不进入冻结画像或
质量排序。V1 formal candidate 只消费 \(S_G\)，旧 \texttt{global\_lcb} 字段必须拒绝。

\paragraph{Contract v5 clarification.}
Strong Global is frozen once, with order \(S_G\rightarrow R_{peer,stable}\rightarrow R_{LOO,medoid}\rightarrow\) frozen random; a tie-break layer is used only when every candidate in the current tie group is evaluable. An unsupported risk/family component, ambiguous family activation, or insufficient conditional support sets only that component to zero. Overall Full fallback is reserved for out-of-support tasks, profile-version conflict, or endpoint ranking instability.

% PAPER_A_METHOD_CONTRACT_CURRENT.json paper_a_method_20260810_v20 SHA-256 085fd9ec0f14a93c986e9f7dda7883173d9ec64959188139a25012d0e6d36b9d
\section{T1：Semi-Auto 条件效应}

\subsection{设计}

\[
2\times2:
\quad
Manual/Semi
\times
ordinary/stress_{\text{assist}}.
\]

50:50 是实验设计分布，不代表生产分布。

T1 使用独立的 T1 gate：\texttt{CALIBRATION\_ENROLLMENT\_CLOSED}、
\texttt{ALL\_CALIBRATION\_WORKERS\_TERMINAL}、\texttt{C1\_EVIDENCE\_FROZEN}、
\texttt{C2\_B\_FROZEN}、\texttt{C2\_A\_RP\_CLOSED}、\texttt{FINAL\_POOLED\_PROFILE\_FROZEN}、
\texttt{T1\_ROSTER\_FROZEN}、\texttt{T1\_TASK\_POOL\_FROZEN}、
\texttt{T1\_RANDOMIZATION\_PLAN\_FROZEN} 和 \texttt{T1\_SAP\_FROZEN}。
T1 不要求 \texttt{STRONG\_GLOBAL\_FROZEN}、\texttt{FULL\_POLICY\_FROZEN} 或
\texttt{VALIDATION\_ROSTER\_FROZEN}，也不以 Global 排名替代本节预先冻结的 2$\times$2 条件设计。
每个 role 必须有可验证的冻结 artifact、SHA 和依赖关系。

\subsection{工人随机分配}

每图：

\begin{verbatim}
2 Manual
2 Semi
\end{verbatim}

四个 slot 在合格工人中约束随机分配：

\begin{itemize}
  \item 同一工人不看同图两种模式；
  \item 每名工人 Manual/Semi 平衡；
  \item ordinary/stress 平衡；
  \item workload cap；
  \item 保存 candidate set、seed 和 assignment probability。
\end{itemize}

\subsection{结构有效性}

每条提交先判：

\begin{verbatim}
structurally_valid
worker_caused_structural_failure
external_system_failure
not_evaluable
\end{verbatim}

worker-caused structural failure 取 \texttt{structurally\_valid=0}，且在其原 Manual/Semi 条件内保留。external system failure 不归责工人或条件；仅在冻结证据规则确认、结果尚不可见时，才触发第 4.6 节的处置。其他无法归因异常为 not evaluable，不静默并入任一失败类别。

\subsection{Primary quality estimand}

对每条提交：

\[
U_{t,u}
=
I(\text{structurally valid})\times IoU(A_{t,u},G_t).
\]

worker-caused structural failure 的 submission-level delivery-adjusted quality 取 0。确认的 external system failure 不进入该提交质量分母：必须重跑同一 image 的完整 Manual/Semi 配对一次，且保持原条件、冻结版本和 worker-image 隔离；无法完成完整配对重跑时，整对行政删失。行政删失不是零值，也不得只删去某一条件的不利结果。

图片--条件层：

\[
\bar U_{t,c}
=
\frac{1}{2}\sum_{u\in c}U_{t,u}.
\]

图片内差：

\[
D_t=\bar U_{t,Semi}-\bar U_{t,Manual}.
\]

这是主要 image-level paired quality estimand。

\texttt{pair\_id} 只用于随机化平衡、事故重跑和 provenance，不是正式独立统计单位。
每图预先冻结两个 pair；pair-level 重跑结束后仍映射回原 image，两个 pair 合并成唯一
image-level outcome。主要推断以 image 为单位，并按 building/image 及重复 worker
结构采用冻结的聚类方法。

每个 T1 条件及其配对结果同时报告原始提交/配对数、worker-caused structural failure 数、external incident 重跑数、行政删失数和原因。所有这些处理在结果不可见时冻结，不能根据质量差异重分类。

\subsection{Secondary outcomes}

\begin{verbatim}
structurally valid rate
valid-only GT IoU
pairwise agreement
worst individual GT quality
blind trust
correction failure
over-correction
Model Issue recognition
\end{verbatim}

两人没有天然多数，因此除非预先冻结双标注融合算法，否则不把“两人聚合 GT IoU”作为正式主结果。

\subsection{Active time}

使用完整 owner-valid active time。

active-time 的确认性效应以 Semi 减 Manual 的条件差
\(\Delta_T=T_{\mathrm{Semi}}-T_{\mathrm{Manual}}\) 报告，负值表示 Semi 更快；
log-ratio 作为偏态分布下的敏感性尺度。只有同时满足总体 coverage 门槛、且
Manual/Semi 与 ordinary/stress 的差异性缺失不超过 SAP 预先冻结的容忍度时，
active time 才进入确认性检验；否则仅作 descriptive 与 missingness sensitivity。

不进行：

\begin{itemize}
  \item 固定减去估算的 Model Issue 时间；
  \item 给 Manual 增加无意义对称表单。
\end{itemize}

\subsection{统计层级}

\begin{enumerate}
  \item delivery-adjusted quality 的非劣，并同时通过 structural-validity safety gate；
  \item owner-valid active time 的预定义差值尺度，log-ratio 为敏感性分析；
  \item mode\(\times\)risk\_assist 交互；
  \item blind trust、correction failure、over-correction 和 Model Issue recognition 等机制。
\end{enumerate}

质量非劣 margin、结构安全门和 active-time 的确认性尺度均来自结果可见前冻结的
Statistical Analysis Plan 与 analysis manifest；正文只报告其预先冻结的方向和判定规则，
不得根据观察到的效应回调 margin。

行为机制不替代前述确认性结果；若 owner-valid active-time coverage 未达到冻结门槛，
active time 降级为 descriptive/sensitivity，不影响质量与结构主分析。

\subsection{不可评价与整图行政删失}

交付状态与失败归因分开保存：\texttt{delivery\_status} 取 valid/invalid/unavailable，
\texttt{failure\_attribution} 取 worker/external/unknown/none。唯一允许的完整 pair 重跑后，
只要任一预冻结 Manual/Semi pair 仍不可评价，整张 image 从主要 image-level paired estimand
行政删失；不可评价不编码为 0，也不使用剩余 pair 改写主要 estimand。可用 pair 只作 sensitivity。

\paragraph{Contract v5 clarification.}
After the single permitted complete-pair rerun, if any pre-frozen Manual/Semi pair for an image remains not evaluable, the entire image is administratively censored from the primary paired estimand. A usable pair may be reported only in the predeclared sensitivity analysis.

% PAPER_A_METHOD_CONTRACT_CURRENT.json paper_a_method_20260807_v19 SHA-256 60e60f01d762ee89d788979ad9ef243db1887f6ef81f4d5bb8ae2e6f29b5af38
\section{V1：共享容量条件下的 Support-aware Full 政策试验}

\subsection{试验臂}

\begin{verbatim}
Strong Global
vs
Support-aware Full
\end{verbatim}

主要使用 Manual。

V1 使用独立的 V1 gate：\texttt{CALIBRATION\_ENROLLMENT\_CLOSED}、
\texttt{ALL\_CALIBRATION\_WORKERS\_TERMINAL}、\texttt{C1\_EVIDENCE\_FROZEN}、
\texttt{C2\_B\_FROZEN}、\texttt{C2\_A\_RP\_CLOSED}、\texttt{FINAL\_POOLED\_PROFILE\_FROZEN}、
\texttt{STRONG\_GLOBAL\_FROZEN}、\texttt{FULL\_POLICY\_FROZEN}、
\texttt{VALIDATION\_ROSTER\_FROZEN} 和 \texttt{V1\_SAP\_FROZEN}。
这些 roles 都必须绑定对应冻结 artifact、SHA 和依赖关系；C2-B closeout 只提供审计证据，
不能直接替代 Full Policy。

\subsection{Block 级共享容量合同}

每个 block 前冻结：

\begin{verbatim}
availability_snapshot
worker_total_capacity
global_quota_per_worker
full_quota_per_worker
candidate_roster
offer_timeout
completion_timeout
max_offer_attempts
replacement_rule
invalid_submission_replacement_rule
candidate_exhaustion_rule
\end{verbatim}

\subsection{两套独立容量账本}

\begin{verbatim}
capacity_global[u]
capacity_full[u]
\end{verbatim}

主分析不允许跨臂借用。

两臂必须共用：

\begin{verbatim}
offer timeout
completion timeout
maximum offer attempts
decline/no-response replacement
invalid submission replacement
candidate exhaustion terminal rule
\end{verbatim}

唯一差异是推荐排序。

\subsection{调度证据}

记录：

\begin{verbatim}
block_id
policy_arm
candidate_set_at_decision
availability_snapshot_id
policy_recommended_worker
recommendation_rank
offered_worker
offer_sequence
accepted_worker
completed_worker
replacement_reason
timeout
capacity_before
capacity_remaining
candidate_exhausted
\end{verbatim}

区分：

\begin{verbatim}
recommended_not_offered
offered_declined
accepted_not_completed
completed_invalid
completed_valid
\end{verbatim}

\subsection{动态追加}

初始与追加规则在两臂完全相同。

可冻结候选：

\begin{verbatim}
k_initial = 2
formal_structure_min_k = 3
standard_cap = 4
exceptional_cap = 5
\end{verbatim}

\(k_{\mathrm{initial}}=2\) 只表示首次 offer 数，不代表可以在两条合法提交时
形成正式 resolved 输出。crowd structure 与 stopping 至少要求三名不同且合资格的
worker IDs；未达到该数目时继续追加，至 cap 仍不足则记为 unresolved。实际 cap 与
其他调度数值由 C1/C2 replay 与功效/预算模拟冻结。

\subsection{GT-blind 聚合}

medoid 选择必须只读取合法提交、
冻结几何相似度、兼容簇、结构有效性和冻结 tie-break；不得读取 GT、公开 GT 评价或
依赖政策分数来决定 GT-blind medoid。

对合法提交：

\begin{enumerate}
  \item 将有冻结证据的 external system failure 移入重跑或行政删失处置，不作为工人或政策质量证据；
  \item 将 worker-caused structural failure 记录为 worker failure event，不进入合法聚合候选；
  \item 使用冻结几何相似度寻找最大内部一致簇；
  \item 判断 largest cluster、second cluster、medoid margin 和结构有效性；
  \item resolved 时选择主簇 medoid；
  \item medoid tie 依次使用 geometry SHA、canonical annotation identity 和
  task-bound frozen seed；这些 tie-break 不读取 Global、Full、GT 或 worker quality；
  \item 稳定多峰不强行输出；
  \item 到上限仍不稳定则 unresolved。
\end{enumerate}

\subsection{运行中间状态与最终政策终态}

\begin{verbatim}
external_system_failure_pending_disposition
\end{verbatim}

最终政策终态仅为：

\begin{verbatim}
resolved
unresolved
severe_failure
\end{verbatim}

\subsubsection{resolved}

产生合法 GT-blind 输出 \(O_t^\pi\)。

\subsubsection{unresolved}

存在合法提交，但未形成稳定单一输出。

\subsubsection{severe\_failure}

到上限仍没有可交付合法输出。

\subsubsection{external\_system\_failure\_pending\_disposition}

外部事故确认后不直接成为可比较的政策终态：按第 3.5 节在盲态下同臂、同冻结版本、对称预留容量重跑一次；不满足任一条件则行政删失。重跑不得跨臂借容量、不得使用不同冻结版本，也不得由 V1 结果触发。

\subsection{V1 outcomes}

\subsubsection{非交付指标}

\[
D_t^\pi
=
I(\text{unresolved or severe failure}).
\]

\subsubsection{Severe failure}

\[
H_t^\pi
=
I(\text{severe failure}).
\]

\subsubsection{Resolved-only quality}

\[
Q_t^\pi=IoU(O_t^\pi,G_t).
\]

\subsubsection{交付调整质量}

\[
U_t^\pi
=
I(\text{resolved})\times IoU(O_t^\pi,G_t).
\]

unresolved/severe failure 的 \(U_t^\pi=0\)，表示没有交付正式布局，不表示真实几何 IoU 必然为 0。

candidate exhaustion、替补规则失败和容量耗尽属于 policy-caused failure：保留在原政策臂 ITT，\texttt{policy\_failure=1}；若任务终态为 unresolved 或 severe failure，则交付调整质量为 0。worker invalid submission 作为 worker failure event 保留；若按冻结替补后任务 resolved，不把最终任务交付调整质量改写为 0。external system failure 的成功重跑以重跑任务保留在原臂 ITT；无法合规重跑的行政删失不进入质量分母，但必须单列原始、重跑、删失和事故原因的两臂分布。

严格集合口径为：randomized set 包含全部随机化任务；primary modified-ITT（mITT）
排除预注册且臂盲确认的 external administrative censor；unresolved、severe failure
和 policy-caused failure 仍留在原臂 mITT，交付调整质量为 0。每臂同时报告
randomized、rerun、administrative censor、mITT、resolved、unresolved 与 severe 数量。

\subsection{检验层级与确认性 margin}

\begin{enumerate}
  \item severe failure 不劣；
  \item unresolved＋severe failure 不劣；
  \item 交付调整政策质量；
  \item 标注数量或成本。
\end{enumerate}

resolved-only output quality、\(k_{used}\)、active time、完成时间和
policy\(\times\)risk\_route interaction 为次级结果。severe failure margin、
unresolved+severe margin、delivery-adjusted quality margin 及成本判定尺度均来自结果
可见前冻结的 Statistical Analysis Plan 与 policy analysis manifest；四级层级和 margin
不得由 V1 结果回调。容量、offer、replacement、dynamic redundancy、GT-blind aggregation
和 ITT 状态机的完整字段合同见附录 F--G。

\subsection{事后审查}

事后人工审查只用于：

\begin{itemize}
  \item 解释 unresolved；
  \item 识别 GT 冲突；
  \item 形成反例；
  \item 判断 OOS。
\end{itemize}

不能用人工修订后的输出替换政策输出重新计算 V1 主效应。

也不得在查看 T1/V1 结果后把 worker 或 policy failure 重分类为 external system failure，或选择性重跑某一臂。事故证据、分类时间和重跑/删失决定是独立审计对象。

\paragraph{Contract v5 clarification.}
The online engine consumes only the frozen \texttt{policy\_candidate\_v2.global\_rank\_S\_G}; it does not recompute peer or LOO ordering. Every start validates the method-contract SHA, policy-manifest SHA, candidate-roster SHA, profile version, and Stage 3 closure. The primary hierarchy is severe-failure non-inferiority, unresolved-plus-severe non-inferiority, delivery-adjusted-quality superiority, then count/cost.

% Source: Paper_A_新版完整论文提纲_vFinal_Draft.md §13
\section{试验分布与生产标准化}

\subsection{分层结果}

分别报告：
\[
V_{\pi,o},\qquad V_{\pi,s}.
\]

\subsection{50:50 设计平均}

\[
V_\pi^{\mathrm{design}}
=
0.5V_{\pi,o}+0.5V_{\pi,s}.
\]

只能解释为 balanced experimental mixture。

\subsection{生产标准化}

若目标生产池比例为 \(p_o,p_s\)：
\[
V_\pi^{\mathrm{prod}}
=
p_oV_{\pi,o}+p_sV_{\pi,s}.
\]

比例来自独立自然任务池，不从 50:50 试验样本估计。

上式只直接适用于 delivery-adjusted quality、non-delivery/severe rate 和平均成本等
无条件指标。resolved-only quality 必须按分子与分母标准化：

\[
Q_{\pi,\mathrm{resolved}}^{prod}
=
\frac{p_o r_o q_o+p_s r_s q_s}{p_o r_o+p_s r_s},
\]

其中 \(r_o,r_s\) 为各层 resolved rate，不能机械加权两个条件性均值。

若没有唯一比例，报告：
\begin{itemize}
  \item \texttt{80:20}
  \item \texttt{60:40}
  \item \texttt{50:50}
  \item \texttt{30:70}
\end{itemize}
情景分析。

% Source: Paper_A_新版完整论文提纲_vFinal_Draft.md §14
\section{V2：外部支持与安全退化审计（可选扩展）}

\subsection{目的}

V2 不属于主要贡献，也不再次比较 Global/Full 因果效应；它只检验：
\begin{itemize}
  \item profile 是否支持；
  \item fallback 是否正确；
  \item supported subset 是否低风险；
  \item unsupported subset 是否富集失败；
  \item 是否存在 false-safe。
\end{itemize}

\subsection{指标}

\begin{itemize}
  \item \texttt{supported coverage}
  \item \texttt{risk among supported}
  \item \texttt{unsupported failure enrichment}
  \item \texttt{false-safe rate}
  \item \texttt{catastrophic supported failure}
  \item \texttt{coverage--risk curve}
  \item \texttt{refresh\_required rate}
\end{itemize}

高 fallback 本身不等于成功。

% PAPER_A_METHOD_CONTRACT_CURRENT.json paper_a_method_20260802_v13 SHA-256 ca9c2788261266744b88e7bbd69d253c509a154e8291ddc0cfc3de771eb4232a
\section{统计分析与功效}

\subsection{C1 后模拟 C2-B}

只模拟以下正式候选设计：
\begin{verbatim}
D8  = 4 common anchors + 4 diverse bridges
D10 = 6 common anchors + 4 diverse bridges
D12 = 6 common anchors + 6 diverse bridges
\end{verbatim}

模拟：
\begin{itemize}
  \item 每人任务数
  \item 共同 anchor 数
  \item 多样化 bridge 数
  \item unique image 数
  \item worker--task graph
  \item \(B_u\) interval width
  \item routing activation rate
  \item fallback rate
\end{itemize}

决定 C2-B 的最终结构。

\subsection{T1 功效}

模拟：
\begin{itemize}
  \item worker/image/building 方差；
  \item structurally valid rate；
  \item delivery-adjusted quality；
  \item Manual/Semi correlation；
  \item active-time missingness；
  \item risk interaction。
\end{itemize}

\subsection{V1 功效}

模拟：
\begin{itemize}
  \item profile activation；
  \item policy divergence（仅作可识别性诊断）；
  \item capacity；
  \item refusal/timeout；
  \item dynamic \(k\)；
  \item unresolved；
  \item delivery-adjusted quality；
  \item policy\(\times\)risk interaction。
\end{itemize}

\subsection{Cross-fitted replay}

按 \texttt{image/base\_task} 划 fold。

评价 fold 不得参与：
\begin{itemize}
  \item Global selection
  \item risk feature selection
  \item weight selection
  \item support threshold
  \item fallback
  \item stopping
\end{itemize}

Replay 用于开发、消融和功效，不替代前瞻 V1。

\subsection{多重终点层级}

\subsubsection{T1 确认性}

\begin{enumerate}
  \item delivery-adjusted quality 非劣，并通过 structural-validity safety gate；
  \item owner-valid active time；
  \item risk interaction。
\end{enumerate}

T1 的质量非劣 margin、结构安全门和 active-time 差值尺度在结果可见前写入冻结的
Statistical Analysis Plan 与 analysis manifest；active time 主尺度为
\(\Delta_T=T_{\mathrm{Semi}}-T_{\mathrm{Manual}}\)，log-ratio 仅作敏感性分析。
总体 active-time coverage 以及 Manual/Semi、ordinary/stress 的差异性缺失必须同时
不超过冻结容忍度，否则 active time 降级为 descriptive/missingness sensitivity。

\subsubsection{V1 确认性}

\begin{enumerate}
  \item severe failure 非劣；
  \item unresolved＋severe failure 非劣；
  \item delivery-adjusted policy quality；
  \item 标注数量或成本。
\end{enumerate}

resolved-only quality、\(k_{used}\)、过程时间和 policy\(\times\)risk interaction 为次级
结果。V1 两个安全门和质量非劣 margin 来自冻结的 Statistical Analysis Plan 与
policy analysis manifest，按上述顺序执行，不在观察到结果后改变。

\subsubsection{RQ2}

以预测关联、效应量、方向、支持范围和预测性能为主，不要求每阶段分别达到传统显著性。

\subsubsection{V2}

可选的外部支持与安全退化审计，不作为主要因果结果。

\subsection{C1 variable-k、rolling enrollment 与学习效应}

C1 所有分母按 task--condition--estimand 的 final unique-worker support 计算，并分别报告 target、
original、authorized、late-entry、outside、deficit 与 pooled excess。W034 的 timing 主分析只使用
sentinel 后 owner-valid 的 17 条授权任务；原始缺失 timing 不插补。W034 profile 必须报告
original-only 与 original+authorized sensitivity。

若 rolling enrollment 未激活，pooled 分析与 original-cohort 分析完全相同；若激活，主分析使用
预冻结规则形成的 pooled cohort，并同时报告 original-only sensitivity、late-entry 数量、版本、
窗口和 task exposure。学习效应只进行 completion order、instruction/interface version 与 provenance
的描述性/敏感性分析，不建立结果可见后选择的复杂主模型。所有选择在 Stage 3 gate 前冻结。

\subsection{可识别跨阶段模型与双 margin}

跨 C1/C2 阶段最终模型使用 building random intercept 与 task-within-building random intercept，
并保留 stage fixed effect；不同时放入吸收 stage 的 task fixed effects。聚类同时报告
\(m^{all}=(n_1-n_2)/k\) 与 \(m^{top2}=(n_1-n_2)/(n_1+n_2)\)，并分别报告
\(k_{\mathrm{LOO,medoid}}\) 和 \(k_{\mathrm{LOO,strict}}\)。

\paragraph{Contract v5 clarification.}
No C1+C2 stage effect is interpreted without a frozen cross-stage anchor or equivalent supported structure. The final model uses building and task-within-building random intercepts with a stage fixed effect, whereas C1-only uses a task fixed effect and no stage effect. The V1 quality criterion is superiority after the two delivery hierarchy gates; it is not an outcome-selected quality margin.

% PAPER_A_METHOD_CONTRACT_CURRENT.json paper_a_method_20260810_v20 SHA-256 085fd9ec0f14a93c986e9f7dda7883173d9ec64959188139a25012d0e6d36b9d
\section{结果章节结构}

\subsection{数据完整性与参考组成}

报告：
\begin{itemize}
  \item 公开原始 GT 数量
  \item 极少量纠错 GT 数量
  \item pending/OOS 数量
  \item 各 estimand eligible 数量
\end{itemize}

同时报告 cohort 与 deviation：
\begin{itemize}
  \item C1 assigned cohort：23；\texttt{administratively excluded}：1（W014）；
  \item \texttt{administratively eligible capability cohort}：最多 22；
  \item active-time analytic cohort：按 owner-valid 行级资格和缺失情况进一步减少，
  不因 W034 接手替补任务而整名排除；
  \item 按 estimand 分别报告 \(n_{\mathrm{workers},Q^{GT}}\)、
  \(n_{\mathrm{workers},R^{peer}}\)、\(n_{\mathrm{workers},F^{struct}}\) 和
  \(n_{\mathrm{workers},\mathrm{StrongGlobal\ eligible}}\)，不使用一个总的 capability
  analytic cohort 替代这些分母；
  \item 原始 assignment manifest 与 authorized reassignment addendum 的数量、时间和 SHA。
\end{itemize}

\subsection{P1 的诊断与预测价值}

报告：
\begin{itemize}
  \item P1\(\rightarrow\)C1；
  \item P1\(\rightarrow\)C2-B；
  \item P1\(\rightarrow\)T1；
  \item failure-family discrepancy workers。
\end{itemize}

\subsection{C1 三轨工人状态}

报告：
\begin{itemize}
  \item task-adjusted GT quality 与 EB 收缩估计；
  \item \texttt{R\_peer\_median}、\texttt{R\_LOO\_medoid} 与 strict LOO 的角色区分；
  \item raw/EB structural failure；
  \item 三者相关与分歧。
\end{itemize}

\subsection{C2 风险韧性与支持}

报告：
\begin{itemize}
  \item 群体 stress effect；
  \item 个体收缩斜率；
  \item C2-A-RP 补测；
  \item eligible/uncertain/fallback。
\end{itemize}

\subsection{T1}

报告：
\begin{itemize}
  \item structurally valid rate；
  \item delivery-adjusted quality；
  \item valid-only GT IoU；
  \item active time；
  \item mode\(\times\)\texttt{risk\_assist}；
  \item blind trust/correction。
\end{itemize}

\subsection{V1}

报告：
\begin{itemize}
  \item policy activation/divergence；
  \item scheduling flow；
  \item resolved/unresolved/severe；
  \item delivery-adjusted quality；
  \item resolved-only quality；
  \item \(k\)/time；
  \item policy\(\times\)\texttt{risk\_route}。
\end{itemize}

\subsection{V2}

可选报告 coverage--risk 与 false-safe，不作为主要结果章节。

\subsection{Failure-family 与反例库审计}

报告：
\begin{itemize}
  \item 各触发来源数量、对应来源分母和富集规则；
  \item candidate、已裁决代表案例与未来候选队列的数量；
  \item 各裁决状态数量与 \texttt{unresolved} 数量；
  \item 代表性 \texttt{confirmed\_failure}、\texttt{confirmed\_ambiguity}、protocol/system issue 案例；
  \item 哪些少量代表案例进入正文、附录或未来回归测试。
\end{itemize}

本节明确声明：反例库不是 prevalence 估计样本，且不会自动导致 GT 修订、画像训练、
C2 阈值或 Full 权重选择。每个 candidate 必须同时报告触发来源、对应分母、候选身份、
裁决状态与 evidence SHA；任一必需 provenance 缺失时，该 candidate 只能进入完整性审计，
不得进入反例富集率或失败率的正式分母。完整 challenge/regression bank 不是本轮必要结果。

\subsection{C1 variable-k 与 enrollment flow}

结果表必须逐 task--condition--estimand 报告 \(k_{target}\)、\(k_{original}\)、\(k_{authorized}\)、
\(k_{late}\)、\(k_{final}\)、outside count、deficit 与 pooled excess，并给出 peer cluster absolute
support/share、cluster margin all 与 cluster margin top2。另行报告 W014 排除、W034 original+17、W001 original+3、
W034 B-004/B-022 outside、W034 sentinel validation 和 timing not-evaluable 数量。

Rolling enrollment 未激活时报告新增人数 0；激活时报告 eligible、admitted、assigned、terminal、
进入 pooled profile 和进入 Validation roster 的人数与 SHA-bound flow，同时并列展示 original-only
和 pooled worker profile、Strong Global/Full activation 及 sensitivity。T1 与 V1 分别公开各自
gate kind、required roles、roster SHA、enrollment registry SHA 及递归依赖 closure；T1 不把
Strong Global、Full 或 Validation roster 当作启动前置条件，V1 则必须显式报告这些依赖。

original cohort 的 administratively eligible 上限为 22；late-entry eligible 另记为
\(N_{late}\)，final pooled cohort 的上界为 \(22+N_{late}\)，再按各 estimand eligibility 减少。

\paragraph{Contract v5 clarification.}
The original assigned cohort is 23 and the original administratively eligible cohort is at most 22. When rolling enrollment is active, late-entry workers are reported as \(N_{late}\), and the pooled pre-estimand upper bound is \(22+N_{late}\); it is not capped at 22.

% Source: Paper_A_新版完整论文提纲_vFinal_Draft.md §17
\section{讨论}

\subsection{P1 为何不仅是准入}

讨论高信息挑战对后续自然任务行为的预测价值，以及 admission 后 range restriction。

\subsection{外部正确性、同行一致性和结构有效性的差异}

解释为什么：
\begin{itemize}
  \item 高 LOO 不一定高 GT
  \item 高 GT 不一定高同行一致性
  \item valid-only IoU 会忽略结构失败
\end{itemize}

\subsection{条件画像的价值与边界}

Full 的价值不在于所有任务都个体化，而在于：
\begin{itemize}
  \item 有支持时有限调整
  \item 无支持时明确 fallback
\end{itemize}

因此，论文主张的是“有支持时有限调整、无支持时安全退化”，而不是个体化必然优于
全局策略。P1 预测、多峰共识和 GT 冲突审查提供解释性与可信度证据，但不替代三项一级贡献。

\subsection{平衡试验与生产分布}

解释 50:50 设计用于识别交互，生产结论需标准化。

\subsection{单一 operational reference 的协议性质}

说明本文只保留一个 operational reference，但不把该协议性选择夸大为不可争议的客观真理。

\subsection{共享工人池中的政策干扰}

讨论对称 quota 和独立容量账本如何定义本文的目标部署制度。

\subsection{局限}

\begin{itemize}
  \item C1 assigned cohort 为 23 人，其中 W014 行政排除后，行政上具备 capability 分析
  资格的 cohort 最多为 22；各 estimand 仍需按 worker-level eligibility 与行级缺失分别
  报告，不能把 W034 因接手替补任务而整名排除；
  \item 条件画像支持有限；
  \item P1 admission 造成范围限制；
  \item 外部域 V2 样本有限；
  \item 部分任务最终 unresolved；
  \item 跨阶段 component 只有在方向、支持和收缩估计均稳定时才进入 Full；不满足时安全
  退化为 component adjustment=0 或 Strong Global；
  \item Full 的增量依赖任务侧 risk activation 的准确性；
  \item 反例库是富集的机制性审计样本，不应被解释为总体失败率；缺失 provenance 的
  candidate 只能保留在完整性审计中。
  \item C1 的 (k) 是 task--condition--estimand-specific；授权替补与 rolling enrollment 改变的
  是可用支持而非原始 assignment target，小 (k) 和 cohort composition 仍限制个体画像精度；
  \item W034 的 17 条授权任务不能修复 sentinel 前的历史 timing missingness，original-only 与
  augmented sensitivity 必须共同解释；
  \item rolling enrollment 若激活会引入 calendar time、instruction exposure 和 completion-order
  差异，因此 pooled 结论必须与 original-only sensitivity 并列，学习效应仅作预先限定的描述。
\end{itemize}

% Source: Paper_A_新版完整论文提纲_vFinal_Draft.md §18
\section{结论}

本文建立了一条从高信息 PreScreen 到前瞻路由验证的完整证据链。核心贡献是三轴工人测量、
双波 Calibration 与部分收缩、以及 Support-aware Full 对 Strong Global 的前瞻政策比较。
P1 形成可检验的 failure-family 假设；C1/C2 分离外部正确性、同行兼容性和结构可靠性，
并以 common anchor、diverse bridge 与精度补测冻结有限条件画像；T1 检验 Semi-Auto 在不同
模型风险下的质量与效率；V1 则在严格控制共享容量、offer、替补、动态冗余和聚合终态后，
比较支持度感知 Full 与 Strong Global。

论文的核心主张不是个体化必然优于全局策略，而是：

\begin{quote}
高信息诊断可以产生具有预测价值的条件假设；当这些假设在独立 Calibration 中得到支持并且任务侧可被可靠激活时，它们能够形成可审计的条件路由；当支持不足时，系统应明确退化为强 Global，而不是强行使用不稳定画像。
\end{quote}


\clearpage
\appendix
\part{完整方法附录}
% Source: Paper_A_新版完整论文提纲_vFinal_Draft.md Appendix A
\section{Reference registry}

\begin{itemize}
  \item 公开原始 GT；
  \item 极少量纠错 GT；
  \item blind preliminary adjudication；
  \item submission-informed revision。
\end{itemize}

% Appendix B: peer compatibility is not external GT quality.
\section{Geometry LOO 与同行兼容性}

\[
S(A_u,A_v)=\min\{S_{\mathrm{boundary}}(A_u,A_v),
S_{\mathrm{wall}}(A_u,A_v)\}.
\]

\begin{align*}
R^{peer}_{t,u}&=\operatorname{median}_{v\neq u}S(A_{t,u},A_{t,v}),\\
R^{peer\_all}_{u}&=\operatorname{median}\{S(A_{t,u},A_{t,v}):v\neq u,\ v\text{ eligible}\},\\
R^{peer\_nonmultimodal}_{u}&=\operatorname{median}\{S(A_{t,u},A_{t,v}):v\neq u,\ v\text{ eligible},\
t\notin\texttt{supported\_multimodal}\},\\
R^{LOO}_{t,u}&=S(A_{t,u},M_t^{(-u)}).
\end{align*}

报告：
\begin{itemize}
  \item \texttt{R\_peer\_all}：全任务描述性同行兼容指标；
  \item \texttt{R\_peer\_nonmultimodal}：Global tie-break 与 stable-compatibility 的正式同行指标；
  \item \texttt{R\_LOO\_medoid}：唯一稳定主簇存在时的高置信参考；
  \item \texttt{R\_LOO\_strict\_stable}：保守敏感性分析；
  \item similarity、medoid、margin、最大/第二簇支持、多峰状态和 structural failure attribution。
\end{itemize}

3:2 任务标记为 \texttt{supported\_multimodal}，不生成唯一 crowd truth；两个簇分别与 GT
比较。LOO 只用于兼容性、聚合稳定性、tie-break 和敏感性分析，不替代外部正确性，也不与
同行 median 机械平均。支持多峰任务不进入 Global 的稳定兼容性证据，仅进入 discrepancy
与敏感性分析。

% Source: Paper_A_新版完整论文提纲_vFinal_Draft.md Appendix C
\section{P1 predictive evidence}

\begin{itemize}
  \item failure-family；
  \item P1\(\rightarrow\)C1/C2/T1；
  \item range restriction；
  \item support。
\end{itemize}

% PAPER_A_METHOD_CONTRACT_CURRENT.json paper_a_method_20260803_v18 SHA-256 694a3126342d7c8de4a5ed788d7ac50b2fec4f104d0b77ace0e604d078c39a87
\section{C2 design simulation}

正式候选仅为：
\begin{verbatim}
D8  = 4 common anchors + 4 diverse bridges
D10 = 6 common anchors + 4 diverse bridges
D12 = 6 common anchors + 6 diverse bridges
\end{verbatim}

模拟并审计：
\begin{itemize}
  \item worker--task 图连通性、common anchor 和 ordinary/stress 支持；
  \item task/building 多样性、Global 排名稳定性和风险斜率精度；
  \item 预算与 policy divergence（仅作 V1 可识别性诊断）；
  \item C2-A-RP 正式只派发 \texttt{risk\_slope\_precision}；top-$k$/component 仅作诊断；
  \item 采用冻结 threshold manifest 的 \texttt{normal\_95\_max\_unified\_slope\_sd} 95\% 正态区间，
  每个 block 为 1 ordinary + 1 stress，最多两个 block、0/2/4 张任务；
  \item 每个 block 后从 canonical、valid、risk-slope-estimand-eligible submissions 重新拟合同一模型，
  区分 declared/projected 与 observed support；达到目标即停，两块后仍不达标时 adjustment=0 与
  Strong Global fallback。
\end{itemize}

C2-S 仅作为可选诊断，不是正式主线阶段。

% PAPER_A_METHOD_CONTRACT_CURRENT.json paper_a_method_20260803_v18 SHA-256 694a3126342d7c8de4a5ed788d7ac50b2fec4f104d0b77ace0e604d078c39a87
\section{Strong Global/Full policy specification}

\subsection{Strong Global}

正式分数为
\[
S^G_u=z(\widehat Q^{GT,EB}_u),
\]
并受 process、independence、GT support、building/task coverage、结构失败安全门和
production quality floor 约束。LOO 只作 tie-break、discrepancy、聚合稳定性和敏感性分析。

\subsection{Support-aware Full}

\[
S^F_{u,t}=S^G_u+I_{\mathrm{risk},t}B_u+I_{f,t}H_{u,f}.
\]

最多两个条件调整：ordinary/stress 风险韧性和一个已激活 failure-family component。
正式 family 仅为 undercoverage/underextension、adjacent-space overextension、
corner/topology instability。每任务最多激活一个 family；总调整有上限；不确定性大时
向群体均值收缩；无支持或超出支持范围时 adjustment=0 或整体 fallback Strong Global。

\subsection{证据门}

component 必须通过 P1 integrity、C1 predictive validation、C2-B confirmation，具备足够
worker/task support、稳定方向、收缩估计和标注前 activation。T1/V1 outcome 不得用于修改
这些规则、阈值或权重。stage-stacked 回归/层级模型必须包含 P1 worker component、
stage、P1\(\times\)stage interaction 以及 task/worker 聚类或随机效应，并输出方向、
stage interaction、worker/task support 和收缩估计。若 C2 出现清晰方向反转、支持不足
或收缩不稳定，component adjustment=0，不得进入正式 Full。

\subsection{Formal freeze requirements}

Strong Global 正式主分数固定为 \(z(\widehat Q^{GT,EB}_u)\)，z 参数使用 gate 前冻结 cohort；
administrative、estimability、reference、process、independence、support、structural EB 或 LCB floor
缺失时 fail closed。Full 必须绑定 risk/family estimate、support、shrinkage、interval、weight、cap、
threshold、margin、profile version、fallback、minimum worker/cluster rule 和 input SHA。T1 assignment
消费 T1 gate；V1 manifest 只消费独立且 SHA 匹配的 V1 gate。相应 gate 未开放、依赖 SHA 变化或
manifest 未声明该 gate 时，禁止产生正式 assignment。

% Source: Paper_A_新版完整论文提纲_vFinal_Draft.md Appendix F
\section{V1 scheduling state machine}

\begin{itemize}
  \item availability；
  \item quota；
  \item offer；
  \item timeout；
  \item replacement；
  \item capacity；
  \item exhaustion。
\end{itemize}

% Source: Paper_A_新版完整论文提纲_vFinal_Draft.md Appendix G
\section{V1 aggregation and terminal state}

\begin{itemize}
  \item valid submission；
  \item largest cluster；
  \item GT-blind medoid；
  \item resolved；
  \item unresolved；
  \item severe/system failure。
\end{itemize}

medoid 只能读取合法提交、冻结几何相似度、兼容簇、结构有效性和冻结 tie-break，
不得读取 GT、Global、Full 或 worker quality。tie-break 依次使用 geometry SHA、
canonical annotation identity 和 task-bound frozen seed；正式 resolved 至少需要三名
不同且合资格的 worker IDs。两名合法提交只构成追加状态，达到 cap 仍不足时记为
\texttt{unresolved}。

% PAPER_A_METHOD_CONTRACT_CURRENT.json paper_a_method_20260803_v18 SHA-256 694a3126342d7c8de4a5ed788d7ac50b2fec4f104d0b77ace0e604d078c39a87
\section{Statistical analysis plan}

\begin{itemize}
  \item T1 primary hierarchy：delivery-adjusted quality 非劣及 structural-validity safety
  gate、owner-valid active time、mode $\times$ risk\_assist interaction；机制结果为次级；
  \item V1 primary hierarchy：severe failure 非劣、unresolved+severe failure 非劣、
  delivery-adjusted policy quality、标注数量或成本；resolved-only quality 和过程机制为次级；
  \item T1/V1 的非劣 margins、结构安全门和 active-time 尺度来自结果可见前冻结的
  Statistical Analysis Plan 与 analysis/policy manifest；不可由观察结果回调；
  \item 以 image/pair、worker/building cluster 和 task-level ITT 为推断单位，行政删失不改写为零；
  \item active-time 必须同时满足总体 coverage 与 Manual/Semi、ordinary/stress 差异性
  缺失的冻结容忍度，否则降级为 descriptive/missingness sensitivity；
  \item 50:50 只解释为 balanced experimental mixture，生产标准化使用独立自然任务池权重；
  \item replay 只用于设计、功效、消融和审计，不替代前瞻 V1；
  \item P1 predictive validation、C2-B confirmation 和 C2-A-RP precision completion 不混用；
  \item T1/V1 outcome 不得回流修改 admission、Calibration、risk、threshold、weight 或 stop rule。
  \item C1 support 按 task--condition--estimand 的 final unique worker 计算；outside 和 duplicate
  revision 不进入分母，peer 先 task-level 后 worker-level 等权汇总；
  \item W034 timing 只使用 sentinel 后 owner-valid 授权行，不回填历史缺失；报告 original-only 与
  original+authorized sensitivity；
  \item rolling enrollment 默认关闭；激活时报告 pooled 主分析与 original-only sensitivity，
  并在 T1 前冻结 enrollment registry、pooled profile 与 T1-specific artifacts，在 V1 前另外
  冻结 Strong Global、Full、Validation roster、V1 SAP 与全部对应 SHA。
\end{itemize}

% PAPER_A_METHOD_CONTRACT_CURRENT.json paper_a_method_20260803_v18 SHA-256 694a3126342d7c8de4a5ed788d7ac50b2fec4f104d0b77ace0e604d078c39a87
\section{Machine-readable evidence schema}

正式 C2-A-RP 仅接受 \texttt{c2a\_rp\_precision\_plan\_v2} 与 \texttt{c2a\_rp\_assignment\_manifest\_v2} 两个 CSV schema；旧 schema 在 closeout 阶段拒绝。

至少包括：
\begin{itemize}
  \item \texttt{task/reference}、\texttt{worker state}、\texttt{Q\_GT\_EB}、\texttt{F\_struct\_EB}；
  \item \texttt{R\_peer\_task}、\texttt{R\_peer\_all}、\texttt{R\_peer\_stable}、
  \texttt{R\_LOO\_medoid}、\texttt{R\_LOO\_strict}；
  \item \texttt{risk}、\texttt{assignment}、\texttt{offer}、\texttt{capacity}、\texttt{response}；
  \item \texttt{aggregation}、\texttt{terminal outcome}、\texttt{analysis inclusion}、\texttt{freeze version}；
  \item \texttt{gt\_score\_computable}、\texttt{gt\_primary\_analysis\_eligible}；
  \item \texttt{candidate\_id}、\texttt{trigger\_source}、\texttt{trigger\_rule\_version}、\texttt{source\_denominator\_definition}；
  \item \texttt{adjudication\_status}、\texttt{adjudicator}、\texttt{adjudicated\_at}、\texttt{evidence\_sha}；
  \item \texttt{future\_use\_status}，但不把反例库字段用于当前画像、阈值或 Full 权重学习；
  \item 反例字段是目标合同，若 provenance、分母或 evidence SHA 不齐全，结果状态必须为
  \texttt{incomplete}，不得写成完整 bank。
  \item \texttt{assignment\_provenance} 四值、\texttt{late\_entry\_calibration\_assignment}、
  \texttt{outside\_primary\_excluded}、original/authorization/distribution SHA；
  \item 每个 estimand 的 \texttt{k\_original}、\texttt{k\_authorized}、\texttt{k\_late}、
  \texttt{k\_final}、support deficit 与 pooled excess；
  \item peer task statistic、equal-task-weight aggregate、\texttt{cluster\_margin\_all}、
  \texttt{cluster\_margin\_top2}、cluster absolute support/share、cluster count 与
  stable/multimodal status；不得生成已废止的 \texttt{normalized\_cluster\_margin}；
  \item \texttt{F\_struct\_raw}、\texttt{F\_struct\_EB}、interval 与 recurrent serious flag；
  \item W034 sentinel manifest SHA、owner-valid status、validation time、log-source SHA 和
  \texttt{timing\_not\_evaluable\_reason}；
  \item rolling activation、enrollment batch/window/version/seed/workload、original/pooled cohort、
  Stage 3 gate dependency SHA、Validation roster SHA 与 formal minimum worker/cluster rule；
  \item C2-A-RP 的 \texttt{target\_component}、\texttt{gap\_reason}、\texttt{formal\_goal}、
  \texttt{target\_ci\_half\_width}、\texttt{declared\_support\_after}、
  \texttt{observed\_support\_after}、\texttt{terminal\_state}、\texttt{fallback\_action} 与
  \texttt{threshold\_manifest\_sha256}；
  \item C2-B deployment 的 \texttt{deployment\_id}、语言、server/project identity、planned
  import/assignment/design/worker-registry SHA；canonical runtime join 使用
  \texttt{(deployment\_id, project\_id, runtime\_task\_id, worker\_id)}，原始 active-time key 保持兼容；
  \item T1 与 V1 的独立 gate kind、required roles、artifact role、artifact SHA 与递归依赖 closure。
\end{itemize}


\clearpage
\nocite{sun2019horizonnet,dumitrache2018crowdtruth,li2017optimalstopping}
\bibliographystyle{elsarticle-num}
\bibliography{refs/references}

\end{document}
